Checked against Image-Line\textquotesingle s official shortcut reference, not from memory. These are for
FL Studio on Windows.

Print this, or keep it open on a second monitor for the first month.

\subsection{Windows}\label{fl04_shortcuts-windows}

{\def\LTcaptype{none} % do not increment counter
\begin{longtable}[]{@{}ll@{}}
\toprule\noalign{}
Key & Action \\
\midrule\noalign{}
\endhead
\bottomrule\noalign{}
\endlastfoot
F5 & Playlist \\
F6 & Channel Rack / Step Sequencer \\
F7 & Piano Roll \\
F8 & Plugin Picker \\
F9 & Mixer \\
F10 & Settings (audio, MIDI, general, file) \\
F11 & Song info \\
F12 & Close all windows \\
Alt+F8 & Sample Browser \\
F1 & Help \\
\end{longtable}
}

\subsection{Transport}\label{fl04_shortcuts-transport}

{\def\LTcaptype{none} % do not increment counter
\begin{longtable}[]{@{}ll@{}}
\toprule\noalign{}
Key & Action \\
\midrule\noalign{}
\endhead
\bottomrule\noalign{}
\endlastfoot
Space & Play / stop \\
Ctrl+Space & Play / pause (does not return to start) \\
L & Toggle Pattern / Song mode \\
R & Toggle recording \\
Ctrl+M & Metronome \\
Ctrl+P & Recording count-in \\
Ctrl+H & Panic, stop all sound \\
\end{longtable}
}

\subsection{Patterns}\label{fl04_shortcuts-patterns}

{\def\LTcaptype{none} % do not increment counter
\begin{longtable}[]{@{}ll@{}}
\toprule\noalign{}
Key & Action \\
\midrule\noalign{}
\endhead
\bottomrule\noalign{}
\endlastfoot
F4 & Create next empty pattern, with naming dialog \\
F2 & Rename the selected pattern \\
\end{longtable}
}

\subsection{Files}\label{fl04_shortcuts-files}

{\def\LTcaptype{none} % do not increment counter
\begin{longtable}[]{@{}ll@{}}
\toprule\noalign{}
Key & Action \\
\midrule\noalign{}
\endhead
\bottomrule\noalign{}
\endlastfoot
Ctrl+S & Save \\
Ctrl+Shift+S & Save as \\
Ctrl+N & \textbf{Save new version.} Not "new file" \\
Ctrl+R & Export WAV \\
Ctrl+Shift+R & Export MP3 \\
Ctrl+Shift+M & Export MIDI \\
\end{longtable}
}

The Ctrl+N behaviour is the one that catches everyone coming from other software. Use
\textbf{File \textgreater{} New} from the menu.

\subsection{Routing}\label{fl04_shortcuts-routing}

{\def\LTcaptype{none} % do not increment counter
\begin{longtable}[]{@{}ll@{}}
\toprule\noalign{}
Key & Action \\
\midrule\noalign{}
\endhead
\bottomrule\noalign{}
\endlastfoot
Ctrl+L & Route selected channels to free mixer tracks \\
\end{longtable}
}

The most useful shortcut in the program. Select your channels in the Channel Rack, press it,
and every one gets its own mixer track ready for EQ and effects.

\subsection{Tools (Playlist and Piano Roll)}\label{fl04_shortcuts-tools-playlist-and-piano-roll}

{\def\LTcaptype{none} % do not increment counter
\begin{longtable}[]{@{}ll@{}}
\toprule\noalign{}
Key & Tool \\
\midrule\noalign{}
\endhead
\bottomrule\noalign{}
\endlastfoot
P & Draw (pencil) \\
B & Paint (brush) \\
D & Delete \\
E & Select \\
Z & Zoom \\
C & Slice \\
T & Mute \\
Y & Playback \\
S & Slip edit (Playlist only) \\
\end{longtable}
}

\subsection{Editing}\label{fl04_shortcuts-editing}

{\def\LTcaptype{none} % do not increment counter
\begin{longtable}[]{@{}ll@{}}
\toprule\noalign{}
Key & Action \\
\midrule\noalign{}
\endhead
\bottomrule\noalign{}
\endlastfoot
Ctrl+B & Duplicate selection to the right \\
Ctrl+C / Ctrl+V / Ctrl+X & Copy / paste / cut \\
Shift+left-click and drag & Clone a clip or pattern \\
Alt (held while dragging) & Bypass snapping \\
Alt+Q & Quantize (Piano Roll) \\
\end{longtable}
}

Ctrl+B is how you turn an 8-bar section into 32 bars in three keystrokes. Learn it early.

\subsection{Mouse conventions worth knowing}\label{fl04_shortcuts-mouse-conventions-worth-knowing}

FL uses the right mouse button far more than other software.

\textbf{Right-click a knob} for a menu with "Type in value", "Create automation clip", "Link to
controller", and "Reset". This is how most automation gets made.

\textbf{Right-click a step} in the step sequencer to turn it off.

\textbf{Right-click a note} in the piano roll to delete it.

\textbf{Right-click a channel name} for options including Piano roll, Rename, Colour, and Delete.

\textbf{Middle-click and drag}, or hold the right mouse button in empty space, to pan around the
playlist.

\textbf{Scroll wheel} over a knob adjusts it in fine increments. Ctrl+scroll is finer still.

\subsection{The three you will use most}\label{fl04_shortcuts-the-three-you-will-use-most}

If you remember nothing else from this file:

\textbf{Ctrl+L} to route to the mixer.
\textbf{Ctrl+B} to duplicate.
\textbf{F4} to make a named pattern.
